\documentclass{article}

% ICLR 2027 官方投稿样式。投稿阶段请勿启用 \iclrfinalcopy。
\usepackage{iclr2027_conference,times}
\input{math_commands.tex}

\usepackage{amsmath}
\usepackage{amssymb}
\usepackage{booktabs}
\usepackage{graphicx}
\usepackage{multirow}
\usepackage{microtype}
\usepackage{hyperref}
\usepackage{url}

% 工作标题，收到正式写作内容后再统一修改。
\title{GroupOpt: A General Construction Paradigm for Neural Combinatorial Optimization}

% 双盲投稿阶段保持匿名。Camera-ready 时再填写真实作者信息。
\author{Anonymous Authors}

% \iclrfinalcopy  % 仅在论文接收并制作 camera-ready 时启用。

\begin{document}

\maketitle

\begin{abstract}
% TODO: 摘要必须为单段；建议依次说明问题、方法、主要结果和结论。
This document is the anonymous ICLR 2027 manuscript template for GroupOpt.
\end{abstract}

\section{Introduction}
\label{sec:introduction}

% TODO: 研究背景、现有自回归构造的局限、GroupOpt 的核心思想、贡献列表。


\section{Related Work}
\label{sec:related-work}

% TODO: 神经组合优化、序列式解码器、多起点/对称性方法、构造策略。


\section{Method}
\label{sec:method}

% TODO: 问题定义、forest construction state、候选 tail、原生 head decoder、
% Tail Selector、合法性约束、训练目标和与宿主模型的接口。


\section{Experiments}
\label{sec:experiments}

% TODO: 实验设置、TSP50 主结果、跨规模结果、跨问题结果、参数量对照、消融。


\section{Conclusion}
\label{sec:conclusion}

% TODO: 总结方法、主要证据、适用范围和局限。



\subsection*{AI use statement}
% ICLR 2027 必需部分，不计入正文页数。提交前必须按实际使用情况改写。
This section will disclose the use of generative AI tools in accordance with
the ICLR 2027 AI Policy for Authors. The authors will review and verify all
AI-assisted work and take responsibility for the final content of the paper.

\subsection*{Reproducibility statement}
% 推荐部分，不计入正文页数。后续应指向正文、附录和匿名代码。
Implementation details, training protocols, evaluation instances, and paired
statistical analyses will be documented in the main paper and appendix. Code
and experiment configurations will be included in the supplementary material.

% 如论文涉及伦理风险，可在此加入不超过一页的 Ethics statement。

\bibliography{references}
\bibliographystyle{iclr2027_conference}

\appendix
\section{Additional Implementation Details}
\label{app:implementation}

% TODO: 完整超参数、伪代码、合法性说明、更多逐种子结果和可视化。



\end{document}
